\documentclass[../../main.tex]{subfiles}% 注意这里的文件路径不能用 ./main.tex ,否则用latexmk编译子文件会报错
\graphicspath{{\subfix{./image/}}{\subfix{../../image/}}} % 指定图片引用路径,后续可以直接使用图片文件名
% 如果图片存放在上级目录,则使用 ../../image/ 作为备用路径

% 例如:
% \begin{figure}[H]
% \centering
% \includegraphics[width=0.3\textwidth]{图.png}
% \caption{}
% \label{figure:图}
% \end{figure}
% 注意:上述\label{}一定要放在\caption{}之后,否则引用图片序号会只会显示??.

\begin{document}

\section{配方法}

\begin{example}
设 \(f \in C^1[0,1]\) 满足 \(f(0) = 0\).
\begin{align*}
\int_0^1 (1-x^2) |f'(x)|^2 \mathrm{d}x \geqslant 2 \int_0^1 |f(x)|^2 \mathrm{d}x.
\end{align*}
\end{example}
\begin{note}
配方法:我们期望产生
\begin{align*}
\int_0^1 p(x) (f'(x) - a(x)f(x))^2 \mathrm{d}x.
\end{align*}
这种结构. 观察原不等式自然会猜测形如
\begin{align*}
\int_0^1 (1-x^2) (f'(x) - a(x)f(x))^2 \mathrm{d}x.
\end{align*}
分部积分求出 \(a\) 即可.令 
\begin{align*}
J = \int_0^1 (1-x^2) (f'(x) - a(x)f(x))^2 \mathrm{d}x.
\end{align*}
将其展开得
\begin{align*}
J =  & \int_0^1 (1-x^2)|f'(x)|^2 \mathrm{d}x - 2\int_0^1 (1-x^2)a(x) f(x)f'(x)\mathrm{d}x + \int_0^1 (1-x^2)a^2(x)|f(x)|^2 \mathrm{d}x.
\end{align*}
对交叉项进行分部积分得
\begin{align*}
- 2\int_0^1 (1-x^2)a(x) f(x)f'(x)\mathrm{d}x
& = -\left[ (1-x^2)a(x) f^2(x) \right]_0^1 + \int_0^1 f^2(x) \frac{\mathrm{d}}{\mathrm{d}x}\left((1-x^2)a(x)\right) \mathrm{d}x  = \int_0^1 f^2(x) \frac{\mathrm{d}}{\mathrm{d}x}\left((1-x^2)a(x)\right) \mathrm{d}x.
\end{align*}
将此结果代回 \(J\) 的展开式中，合并同类项得
\begin{align*}
J = \int_0^1 (1-x^2)|f'(x)|^2 \mathrm{d}x + \int_0^1 f^2(x) \left[ (1-x^2)a^2(x) + \frac{\mathrm{d}}{\mathrm{d}x}\left((1-x^2)a(x)\right) \right] \mathrm{d}x.
\end{align*}
为了使得拼凑后的结果与我们期望得到的原不等式匹配,即希望方括号内系数为 \(-2\)，我们需要求解如下微分方程：
\begin{align*}
(1-x^2)a^2(x)+\frac{\mathrm{d}}{\mathrm{d}x}\left( (1-x^2)a(x) \right) =-2&\Longleftrightarrow \left( 1-x^2 \right) a^2\left( x \right) -2xa\left( x \right) +\left( 1-x^2 \right) a' \left( x \right) =-2
\\
&\Longleftrightarrow a'\left( x \right) =-a^2\left( x \right) +\frac{2x}{1-x^2}a\left( x \right) -\frac{2}{1-x^2}.
\end{align*}
这是一个Riccati方程,必须先猜出一个特解$y_1(x)$,然后令$a(x)=y_1(x)+\frac{1}{u(x)}$,代入原方程,就能将原方程转化为关于$u(x)$的一阶线性常微分方程,进而得到解.因为这里我们只需要一个解，所以只用猜到一个解就行了，不用解这个方程.观察可知$a(x)=\frac{1}{x}$是一个解.
\end{note}
\begin{proof}
利用配方法, 考虑积分
\begin{align*}
J \triangleq \int_0^1 (1-x^2) \left( f'(x) - \frac{f(x)}{x} \right)^2 \mathrm{d}x.
\end{align*}
展开被积函数可得
\begin{align*}
J = \int_0^1 (1-x^2)|f'(x)|^2 \mathrm{d}x - 2 \int_0^1 \frac{(1-x^2) f(x) f'(x)}{x} \mathrm{d}x + \int_0^1 \frac{(1-x^2) |f(x)|^2}{x^2} \mathrm{d}x 
\end{align*}
对中间项分部积分, 并注意到极限 \(\lim\limits_{x \to 0^+} \frac{f^2(x)}{x} = 0\) 及 \(\lim\limits_{x \to 1^-} \frac{1-x^2}{x} f^2(x) = 0\),可得
\begin{align*}
-2 \int_0^1 \frac{(1-x^2) f(x) f'(x)}{x} \mathrm{d}x  =  \left[ -\frac{1-x^2}{x} f^2(x) \right]_0^1 + \int_0^1 f^2(x) \mathrm{d}\left( -\frac{1-x^2}{x} \right) 
=  \int_0^1 f^2(x) \frac{-x^2-1}{x^2} \mathrm{d}x.
\end{align*}
代回 \(J\) 的展开式得
\begin{align*}
J =   \int_0^1 (1-x^2)|f'(x)|^2 \mathrm{d}x + \int_0^1 f^2(x) \frac{-x^2-1}{x^2} \mathrm{d}x + \int_0^1 \frac{(1-x^2) |f(x)|^2}{x^2} \mathrm{d}x 
=   \int_0^1 (1-x^2)|f'(x)|^2 \mathrm{d}x - 2 \int_0^1 |f(x)|^2 \mathrm{d}x.
\end{align*}
由于 \(J = \int_0^1 (1-x^2) \left( f'(x) - \frac{f(x)}{x} \right)^2 \mathrm{d}x \geqslant 0\), 故原不等式成立.
\end{proof}

\begin{example}
设 \(n \in \mathbb{N}, f \in C^1[0,1]\) 满足 \(f(1) = 0\). 证明:
\begin{align*}
\int_0^1 x^n |f'(x)|^2 \mathrm{d}x \geqslant n \int_0^1 \frac{x^{n-1}}{1-x} |f(x)|^2 \mathrm{d}x.
\end{align*}
\end{example}
\begin{proof}
利用配方法, 考虑积分
\begin{align*}
I \triangleq \int_0^1 x^n \left( f'(x) + \frac{f(x)}{1-x} \right)^2 \mathrm{d}x.
\end{align*}
展开被积函数可得
\begin{align*}
I = & \int_0^1 x^n |f'(x)|^2 \mathrm{d}x + 2 \int_0^1 x^n \frac{f(x)f'(x)}{1-x} \mathrm{d}x + \int_0^1 x^n \frac{|f(x)|^2}{(1-x)^2} \mathrm{d}x \\
= {} & \int_0^1 x^n |f'(x)|^2 \mathrm{d}x + \int_0^1 \frac{x^n}{1-x} \mathrm{d}(f^2(x)) + \int_0^1 \frac{x^n}{(1-x)^2} |f(x)|^2 \mathrm{d}x.
\end{align*}
对中间项分部积分, 并注意到极限 \(\lim\limits_{x \to 1^-} \frac{x^n}{1-x} f^2(x) = 0\),可得
\begin{align*}
2 \int_0^1 x^n \frac{f(x)f'(x)}{1-x} \mathrm{d}x= \left[ \frac{x^n}{1-x} f^2(x) \right]_0^1 - \int_0^1 f^2(x) \mathrm{d}\left( \frac{x^n}{1-x} \right) 
= - \int_0^1 \left( \frac{n x^{n-1}(1-x) + x^n}{(1-x)^2} \right) |f(x)|^2 \mathrm{d}x.
\end{align*}
代回 \(I\) 的展开式得
\begin{align*}
I = {} & \int_0^1 x^n |f'(x)|^2 \mathrm{d}x - \int_0^1 \frac{n x^{n-1}(1-x) + x^n}{(1-x)^2} |f(x)|^2 \mathrm{d}x + \int_0^1 \frac{x^n}{(1-x)^2} |f(x)|^2 \mathrm{d}x \\
= {} & \int_0^1 x^n |f'(x)|^2 \mathrm{d}x - n \int_0^1 \frac{x^{n-1}}{1-x} |f(x)|^2 \mathrm{d}x.
\end{align*}
由于 \(I = \int_0^1 x^n \left( f'(x) + \frac{f(x)}{1-x} \right)^2 \mathrm{d}x \geqslant 0\), 故原不等式成立.
\end{proof}

\begin{example}
设 \(f \in C^1[0,1]\) 满足 \(f(0) = f(1) = 0\), 证明:
\begin{align*}
\int_0^1 x^2 |f'(x)|^2 \mathrm{d}x \geqslant \frac{1}{4} \int_0^1 |f(x)|^2 \mathrm{d}x.
\end{align*}
\end{example}
\begin{proof}
同样利用配方法, 考虑积分
\begin{align*}
J \triangleq \int_0^1 x^2 \left( f'(x) + \frac{f(x)}{2x} \right)^2 \mathrm{d}x.
\end{align*}
展开被积函数可得
\begin{align*}
J = {} & \int_0^1 x^2 |f'(x)|^2 \mathrm{d}x + \int_0^1 x f(x) f'(x) \mathrm{d}x + \int_0^1 \frac{1}{4} |f(x)|^2 \mathrm{d}x.
\end{align*}
对中间项整理并用分部积分:
\begin{align*}
\int_0^1 x f(x) f'(x) \mathrm{d}x  = \frac{1}{2} \left[ x f^2(x) \right]_0^1 - \frac{1}{2} \int_0^1 |f(x)|^2 \mathrm{d}x=- \frac{1}{2} \int_0^1 |f(x)|^2 \mathrm{d}x.
\end{align*}
从而
\begin{align*}
J = {}  \int_0^1 x^2 |f'(x)|^2 \mathrm{d}x - \frac{1}{2} \int_0^1 |f(x)|^2 \mathrm{d}x + \frac{1}{4} \int_0^1 |f(x)|^2 \mathrm{d}x 
= \int_0^1 x^2 |f'(x)|^2 \mathrm{d}x - \frac{1}{4} \int_0^1 |f(x)|^2 \mathrm{d}x.
\end{align*}
由于 \(J = \int_0^1 x^2 \left( f'(x) + \frac{f(x)}{2x} \right)^2 \mathrm{d}x \geqslant 0\), 故原不等式成立.
\end{proof}



\end{document}